%%%PAGE 65 rot=0 legibility=good summary=牛顿运动定律的基本运用与综合运用提纲（两类基本问题、超重失重、瞬时加速度、图像、连接体、板块、传送带、临界极值）
%%%BLOCK id=065-01 ch=03 sec=动力学两类基本问题 kind=summary alt=none cont=none y=0.04-0.15
\textbf{牛顿运动定律的基本运用}
\begin{itemize}
\item 动力学两类基本问题
  \begin{itemize}
  \item 已知受力情况求物体运动状态.
  \item 已知运动情况求物体受力情况\quad 以 $F_{\text{合}}=ma$ 为桥梁,
  \end{itemize}
\end{itemize}
%%%END
%%%BLOCK id=065-02 ch=03 sec=超重失重 kind=summary alt=none cont=none y=0.15-0.22
\begin{itemize}
\item 超重与失重\quad 向上超\ 下失
  \begin{itemize}
  \item 超加失减 $ma$.
  \item 太空人完全失重
  \end{itemize}
\end{itemize}
%%%END
%%%BLOCK id=065-03 ch=03 sec=瞬时加速度 kind=summary alt=none cont=none y=0.22-0.28
\textbf{瞬时加速度问题}
\begin{itemize}
\item 绳杆面、力\unclear{可}瞬间消失\quad 用牛二求 $a_{\text{瞬}}$
\item 弹簧、蹦床、橡皮\unclear{绳}、力保持
\end{itemize}
%%%END
%%%BLOCK id=065-04 ch=03 sec=超重失重 kind=note alt=none cont=none y=0.28-0.32
对超重与失重的理解.
%%%END
%%%BLOCK id=065-05 ch=03 sec=动力学两类基本问题 kind=note alt=none cont=none y=0.34-0.38
动力学的两类基本问题.
%%%END
%%%BLOCK id=065-06 ch=03 sec=动力学图像 kind=summary alt=none cont=none y=0.39-0.46
\textbf{动力学图像问题}
\begin{itemize}
\item $v$-$t$\quad $a$-$t$\quad $F$-$t$\quad $F$-$a$图.\quad 待定系数\ 或\ 利用数据
\end{itemize}
%%%END
%%%BLOCK id=065-07 ch=03 sec=斜面 kind=note alt=none cont=none y=0.47-0.50
三类光滑斜面上下滑时间比较.
%%%END
%%%BLOCK id=065-08 ch=03 sec=连接体与整体隔离 kind=summary alt=none cont=none y=0.52-0.66
\textbf{牛顿运动定律综合运用}
\begin{itemize}
\item 连接体问题
  \begin{itemize}
  \item 整体法.
  \item 隔离法
  \end{itemize}
\end{itemize}
%%%END
%%%BLOCK id=065-09 ch=03 sec=板块模型 kind=method alt=none cont=none y=0.67-0.78
\begin{itemize}
\item 滑块板块模型
  \begin{itemize}
  \item $F$、$a$、$v_{\text{共}}$、$t$、$x$、$Q$、快带慢、快后慢减 % 原稿“快带慢”与“快后慢减”之间上方有一小字 f（疑为摩擦力 f 的批注）
  \end{itemize}
\end{itemize}
%%%END
%%%BLOCK id=065-10 ch=03 sec=传送带 kind=method alt=none cont=none y=0.79-0.87
\begin{itemize}
\item 传送带模型
  \begin{itemize}
  \item 共速临界\quad 摩擦突变
  \end{itemize}
\end{itemize}
%%%END
%%%BLOCK id=065-11 ch=03 sec=临界极值 kind=summary alt=none cont=none y=0.89-1.00
\textbf{动力学中临界极值问题}
\begin{itemize}
\item 分离临界\quad $N=0$\quad $a_1=a_2$\quad $v_1=v_2$
\item 速度临界：合外力\unclear{为0}、收尾速度 % CHECK: “为0”二字连写形似“加”，按物理意义（收尾速度时合外力为零）读作“为0”
\item 滑动临界\quad 静摩擦达最大
\item 绳子松断临界：松 $T=0$；断 $T=\max$
\item \unclear{…} % 页底被照片边缘裁切，最后一行仅露出字的上端，无法辨认
\end{itemize}
%%%END
%%%PAGEEND 65
